\documentclass[10pt,a4paper]{amsart}
\usepackage[margin=19mm]{geometry}
\usepackage{amsmath,amssymb,mathtools}
\usepackage[scr=rsfs,cal=euler]{mathalfa}
\usepackage{xfrac}
\usepackage[unicode,hidelinks,bookmarks=false]{hyperref}
\newcommand{\ie}{i.e.\ }
\newcommand{\cf}{cf.\ }
\newcommand{\resp}{respectively\ }
\newcommand{\Subsection}{\subsection}
\providecommand{\nobreakdash}{\nobreak\hbox{-}}
\setlength{\emergencystretch}{2em}
\multlinegap0pt
\usepackage{latexsym}
\usepackage{ifthen}
\renewcommand{\baselinestretch}{1.2}
\def\eq#1{{\rm(\ref{E#1})}}
\def\Eq#1#2{\ifthenelse{\equal{#1}{*}}
  {\begin{equation*}\begin{aligned}#2\end{aligned}\end{equation*}}
  {\begin{equation}\begin{aligned}\label{E#1}#2\end{aligned}\end{equation}}}
\newtheorem{thm}{Theorem}%[section]
\newtheorem{prop}{Proposition}%[section]
\newtheorem{coro}{Corollary}%[section] 
\newtheorem{lemm}{Lemma}%[section]
\newtheorem{remark}{Remark}%[section]
\newtheorem{exmp}{Example}%[section] 
\newtheorem{claim}{Claim} 
\newtheorem{prclm}{Proof of Claim}[section]
\newtheorem{defin}{Definition}%[section]
\newtheorem{notation}{Notation}%[section]
\newtheorem{probl}{Problem}
\newtheorem{opprob}{Open Problem}
\newcommand{\dltu}{\mathcal{U}}
\newcommand{\dltj}{\mathbf{J}}
\newcommand{\doltj}{\mathcal{J}}
\newcommand{\dltv}{\mathcal{V}}
\newcommand{\dltp}{\mathcal{P}}
\newcommand{\dltw}{\mathcal{W}}
\newcommand{\dlta}{\mathcal{A}}
\newcommand{\dltb}{\mathcal{B}}
\newcommand{\dltc}{\mathcal{C}}
\newcommand{\dltd}{\mathcal{D}}
\newcommand{\dlti}{\mathcal{I}} 
\newcommand{\dlth}{\mathcal{H}}
\newcommand{\vct}[1]{\mathbf{#1}}
\newcommand{\NN}{\mathbb{N}} 
\newcommand{\ZZ}{\mathbb{Z}}
\newcommand{\QQ}{\mathbb{Q}}
\newcommand{\RR}{\mathbb{R}}
\newcommand{\RF}{\mathbb{F}}
\newcommand{\DD}{\mathbb{D}} 
\newcommand{\R}{\mathbf{R}}
\newcommand{\Rn}{\mathbb{R}^n}
\newcommand{\RRp}{\mathbb{R}_+}
\newcommand{\RRpn}{\mathbb{R}_+ ^n}
\newcommand{\RRx}{\overline{\mathbb{R}}}
\newcommand{\epsi}{\varepsilon}
\newcommand{\feqs}{(1) - (n)}
\newcommand{\indhy}{induction hypothesis}
\newcommand{\exmo}{extended monotone left-inverse }
\newcommand{\nma}{\Vert}
\newcommand{\bbs}{\langle}
\newcommand{\jbs}{\rangle}
\newcommand{\bfer}{\Big \lceil}
\newcommand{\jfer}{\Big \rceil}
\newcommand{\baer}{\Big \lfloor}
\newcommand{\jaer}{\Big \rfloor}
\newcommand{\itll}{\mathit{l}}
\newcommand{\chf}{\mathbb{I}}
\newcommand{\seg}{\mbox{\rm s}}
\newcommand{\lne}{\mathit{l}}
\newcommand{\bel}{\mbox{\rm int} \, \mbox{\rm s}}
\newcommand{\card}{\mbox{\rm card}}
\newcommand{\dist}{\mbox{\rm dist}} 
\newcommand{\supp}{\mbox{\rm supp}}
\newcommand{\aff}{\mbox{\rm Aff}} 
\newcommand{\lin}{\mbox{\rm Lin}}
\newcommand{\tria}{\triangle}
\newcommand{\conv}{\mbox{\rm conv}}
\newcommand{\intconv}{\mbox{\rm conv}^{\circ}}
\newcommand{\intseg}{\mbox{\rm s}^{\circ}}
\newcommand{\inttri}{\triangle^{\circ}}
\newcommand{\affdim}{\mbox{\rm affdim}} 
\newcommand{\word}[1]{\text{ \ \ #1 \ \ }}
\newcommand{\eps}{\varepsilon}
\newcommand{\map}{\longrightarrow}
\newcommand{\forevery}{\word{for every}}
\newcommand{\abs}[1]{\vert #1 \vert}
\newcommand{\Abs}[1]{\left\vert #1 \right\vert}
\newcommand{\bsz}[2]{\left \langle #1 \,, #2 \right \rangle}
\begin{document}
\title[Generalized inverses of strictly monotone transformations]{Generalized inverses of strictly monotone transformations}
\author{Péter Tóth}
\date{}
\maketitle
\noindent\emph{Source and licence.} Generalized inverses of strictly monotone transformations. Version arXiv:2608.19997v1 (CC BY 4.0). This material is distributed under the Creative Commons Attribution 4.0 International licence, http://creativecommons.org/licenses/by/4.0/, as recorded in the arXiv OAI licence metadata for this exact version and shown on the abstract page https://arxiv.org/abs/2608.19997v1. Full rights notice: licenses/arxiv-2608.19997-CC-BY-4.0.txt.\par\smallskip
\noindent\textit{Editorial scope.} Counted unit: the original \texttt{thm} environment \texttt{thm\_Appl\_vvQAMs} at source lines 1509--1539 together with its complete proof at lines 1541--1638; the delivered document keeps source lines 1051--1639 so that every referenced label and every citation resolves inside it. Native editable source is the paper's own concatenated TeX; the AMS wrapper is editorial formatting only. No publisher class, upstream script or unrelated artwork is redistributed.
\begin{prop}\label{prop-AngleLimit}
Let $ (x_k) : \NN \map \Rn $ be a sequence, 
$ h \in \Rn $ be a vector and $ \delta > 0 $ be a number 
such that $ \nma h \nma = 1 $, 
$ \nma x_k \nma > \delta $ for all $ k \in \NN $, moreover 
\[
\alpha_k := 
\bsz	{\frac{x_k}{\nma x_k \nma}}{h} \to 1 
\ \mbox{ as } \ k \to \infty. 
\]
Then, for any $ d \in \left( 0, \delta \right) $, 
we have 
\[
\lim_{k \to \infty} 
\bsz{\frac{x_k-dh}{\nma x_k - dh \nma}}{h} 
= 1.
%\to 1. 
\] 
\end{prop}

\begin{proof}
Consider the function $ g(t) := t-2d+\frac{d^2}{t} = 
\frac{1}{t} \left( t-d \right)^2 $ 
defined for all $ t > 0 $. 
Clearly $ g $ is continuous, 
non-negative and 
$ \lim_{t \to \infty} g(t) = + \infty $. 
Since $ \delta > d $, there exists $ M > 0 $ such that 
$ g(t) > M $ for all $ t \geq \delta $. 
In particular, $ g \left( \nma x_k \nma \right) > M $ 
for every index $ k \in \NN $. 
 
Since $ \alpha_k \leq 1 $ due to the CBS inequality, 
we have 
\[ 
\nma x_k \nma - 2d \alpha_k + \frac{d^2}{\nma x_k \nma} \geq 
\nma x_k \nma - 2d + \frac{d^2}{\nma x_k \nma} = 
g \left( \nma x_k \nma \right) > M 
\hspace{10mm} \left( k \in \NN \right). 
\] 
Now we shall calculate 
\begin{align*}
\left( \frac{\nma x_k \nma - d}{\nma x_k - d h \nma} 
\right)^2 
& = 
\frac{\nma x_k \nma ^2 - 2d \nma x_k \nma + d^2}{
\nma x_k \nma ^2 - 2d \bsz{x_k}{h} + d^2} 
= 
1 + 2d \frac{\bsz{x_k}{h}-\nma x_k \nma}{
\nma x_k \nma ^2 - 2d \bsz{x_k}{h} + d^2} \\ 
& = 
1 + 2d \frac{\alpha_k - 1}{
\nma x_k \nma - 2d \alpha_k + \frac{d^2}{\nma x_k \nma}} 
\to 1 
\ \mbox{ as } k \to \infty, 
\end{align*}
because the numerator tends to $ 0 $ while the 
denominator is bounded from below by $ M > 0 $. 
Consequently, the positive sequence 
$ \frac{\nma x_k \nma - d}{\nma x_k - d h \nma} $ 
also tends to $ 1 $. Finally, 
\begin{align*}
1 & \geq \bsz{\frac{x_k-dh}{\nma x_k - dh \nma}}{h} = 
\frac{\nma x_k \nma}{\nma x_k - dh \nma} \cdot 
\bsz	{\frac{x_k}{\nma x_k \nma}}{h} - 
\frac{d}{\nma x_k - dh \nma} = 
\frac{\nma x_k \nma - d}{\nma x_k - d h \nma} \cdot \alpha_k 
\\ & + \frac{d}{\nma x_k - dh \nma} \left( \alpha_k - 1 \right) 
> 
\frac{\nma x_k \nma - d}{\nma x_k - d h \nma} \cdot \alpha_k 
+ \frac{d}{\delta-d} \left( \alpha_k - 1 \right) 
\to 1 \cdot 1 + \frac{d}{\delta-d} \cdot 0 = 1. 
\end{align*}
Therefore 
$ \bsz{\frac{x_k-dh}{\nma x_k - dh \nma}}{h} \to 1 $ 
and the proof is complete. 
%applying the results obtained in the previous steps 
%of the reasoning. Thus the proof is complete.  
\end{proof}

\begin{thm}\label{thm-ExmoContinuous}
Let $ \emptyset \neq C \subseteq \Rn $ be a closed, 
convex set and let $ f : C \map \Rn $ be 
a strictly increasing mapping. Then the unique 
\exmo function 
$ f^{(-1)} : \conv \left( f(C) \right) \map C $ is 
continuous. 
\end{thm}

\begin{proof}
We have to show that if a sequence 
$ \left( u_k \right) : \NN \map \conv (f(C)) $ 
converges to $ v \in \conv (f(C)) $ then 
$ f^{(-1)}(u_k) $ converges to $ f^{(-1)} (v) $. 

Suppose that, on the contrary, 
$ \lim_{k \to \infty} f^{(-1)}(u_k) \neq f (v) $. 
This means that there exists $ \delta > 0 $ 
and a subsequence $ \left( u_{\ell_k} \right) $ of 
$ (u_k) $ such that 
$ \nma f^{(-1)} \left( u_{\ell_k} \right) 
- f^{(-1)}  \left( v \right) \nma > \delta $. 
Without loss of generality we may assume that this 
sequence is the whole $ \left( u_k \right) $. 
Let us observe that 
\[
\frac{f^{(-1)}\left( u_k \right) - 
f^{(-1)}\left( v \right)}{\nma f^{(-1)}\left( u_k \right) - 
f^{(-1)}\left( v \right) \nma} 
\in \mathcal{S}_{n-1} 
\ \mbox{ for all } k \in \NN. 
\]
Since $ \mathcal{S}_{n-1} $ is compact, the above sequence 
has a convergent subsequence. 
That is, there exist $ h \in \Rn $ with 
$ \nma h \nma = 1 $, and strictly 
increasing sequence of indices 
$ (m_k) : \NN \map \NN $ such that if 
$ x_k := 
f^{(-1)}\left( u_{m_k} \right) - 
f^{(-1)}\left( v \right) $ 
for all $ k \in \NN $, then 
\[ 
\left\nma \frac{x_k}{\nma x_k \nma} - h \right\nma 
\to 0 
\ \mbox{ or, equivalently, } \ 
\bsz{\frac{x_k}{\nma x_k \nma}}{h} \to 1 . 
\]
As $ C $ is convex, $ f^{(-1)}(v) + 
\frac{\delta}{\nma x_k \nma} x_k  \in C $ for 
each $ k \in \NN $. But $ C $ is closed as well, so 
$ f^{(-1)}(v) + \delta h \in C $, whence 
\[ 
s := f^{(-1)}(v) + \frac{\delta}{2} h \in C 
\] also holds. 
Now let us apply Proposition \ref{prop-AngleLimit} 
to the sequence $ (x_k) $, the vector $ h $ and 
the constant $ \delta $. The Proposition claims that, 
in particular, for $ d=\frac{\delta}{2} $ the assertion 
\begin{equation}\label{eq-thm_Cont1}
1 = \lim_{k \to \infty} 
\bsz{\frac{x_k-\frac{\delta}{2}h}{
\nma x_k - \frac{\delta}{2}h \nma}}{h} 
= 
\lim_{k \to \infty} 
\bsz{\frac{ f^{(-1)}\left( u_{m_k} \right) - s}{
\nma f^{(-1)}\left( u_{m_k} \right) - s \nma}}{h} 
\end{equation} 
is valid. 
In the next step let us fix another vector 
$ t := \frac{f^{(-1)}(v) + s}{2} \in C $. Then 
\begin{align*}
\frac{\delta}{2} \bsz{h}{f(s)-v} & = 
\bsz{s - f^{(-1)}(v)}{f(s)-v} = 
\bsz{s-t}{f(s)-f(t)} + \bsz{s-t}{f(t)-v} \\ 
& + \bsz{t-f^{(-1)}(v)}{f(s)-v} = 
\bsz{s-t}{f(s)-f(t)} \\ 
& + \bsz{t - f^{(-1)}(v)}{f(t)-v} + 
\frac{1}{2} \bsz{s - f^{(-1)}(v)}{f(s)-v} > 0, 
\end{align*} 
because the first term is positive as $ f $ is strictly 
increasing, while the other two terms are 
non-negative, due to the defining properties of 
$ f^{(-1)} $. Thus 
$ \bsz{h}{f(s)-v} > 0 $, which implies that there 
exists $ \varepsilon > 0 $ and $ k_0 \in \NN $ such that 
\[
\bsz{h}{f(s)-u_{m_k}} > \varepsilon 
\ \mbox{ for all } k > k_0 \, 
\]
since $ u_{m_k} \to v $. 
Furthermore, from \eqref{eq-thm_Cont1} we get that  
\[ 
\left \vert 
\bsz	{f(s) - u_{m_k}}{\frac{ f^{(-1)}\left( u_{m_k} \right) - s}{
\nma f^{(-1)}\left( u_{m_k} \right) - s \nma} - h} 
\right \vert 
\leq 
\left \nma f(s) - u_{m_k} \right \nma \cdot 
\left \nma \frac{ f^{(-1)}\left( u_{m_k} \right) - s}{
\nma f^{(-1)}\left( u_{m_k} \right) - s \nma} - h \right \nma 
\to 0, 
\]
applying the CBS inequality and 
using that $ \left( f(s) - u_{m_k} \right) $ 
is a bounded sequence. 
In particular, there exists an index $ k_1 > k_0 $ 
such that 
\[
\bsz	{f(s) - u_{m_k}}{\frac{ f^{(-1)}\left( u_{m_k} \right) - s}{
\nma f^{(-1)}\left( u_{m_k} \right) - s \nma} - h} 
> - \frac{\varepsilon}{2} 
\ \mbox{ for all } k > k_1. 
\]
Combining the previous results we obtain that, 
for all $ k > k_1 \,$, the inequality 
\begin{align*}
& \bsz{f(s)-u_{m_k}}{\frac{ f^{(-1)}\left( u_{m_k} \right) - s}{
\nma f^{(-1)}\left( u_{m_k} \right) - s \nma}} 
= 
\bsz{f(s)-u_{m_k}}{\frac{ f^{(-1)}\left( u_{m_k} \right) - s}{
\nma f^{(-1)}\left( u_{m_k} \right) - s \nma}-h} \\ 
& + \bsz{f(s)-u_{m_k}}{h} > - \frac{\varepsilon}{2} 
+ \varepsilon = \frac{\varepsilon}{2} > 0 
\end{align*} 
holds. 
But this implies 
\[
\bsz{f^{(-1)}\left( u_{m_k} \right) - s}{u_{m_k} - f(s)} 
< 0 
\hspace{10mm} \left( k > k_1 \right), 
\]
which contradicts the monotonicity of $ f^{(-1)} $. 
Therefore our initial assumption about 
$ f^{(-1)}(u_k) $ not converging to $ f^{(-1)}(v) $ 
was false, so $ f^{(-1)} $ is continuous. 
\end{proof}

\section{Application for quasi-arithmetic means}\label{Section-Appl} 

In the final section of our paper we introduce the 
concept of vector valued quasi-arithmetic means. 
\begin{defin}
Let $ \emptyset \neq C \subseteq \Rn $ be a 
closed, convex set and let $ f : C \map \Rn $ 
be a strictly monotone mapping. Moreover, 
let $ m \in \NN $ and $ \lambda_1 \,, \dots , \lambda_m > 0 $ 
be fixed numbers such that $ m \geq 2 $ and 
$ \lambda_1 + \dots + \lambda_m = 1 $. 

Then the function 
$ \mathcal{M}_{f, \lambda_1 , \dots , \lambda_m}^{[m]} 
: C^m \map C $ defined by 
\[
\mathcal{M}_{f, \lambda_1 , \dots , \lambda_m}^{[m]} 
\left( x_1 \,, \dots , x_m \right) 
:= f^{(-1)} \bigl( \lambda_1 f(x_1) + \dots 
+ \lambda_m f(x_m) \bigr) 
\hspace{8mm} 
\left( x_1 \,, \dots , x_m \in C \right) 
\] 
is called the 
{\em $ m $-variable weighted quasi-arithmetic mean}. 
The {\em generator} of the mean is $ f $. 
\end{defin}

It is easy to see that for arbitrary weights 
$ \mathcal{M}_{f, \lambda_1 , \dots , \lambda_m}^{[m]} = 
\mathcal{M}_{-f, \lambda_1 , \dots , \lambda_m}^{[m]} \, $, 
thus it is enough to consider the case of a strictly 
increasing generator. In \cite{Leo23} the idea of 
vector-valued QAMs appears, although it is 
{\em a priori assumed} that the generator has a convex image, 
which we find rather restrictive. 

One of the first natural problems which arises 
immediately is the equality problem: for a 
fixed number of variables and fixed weights find 
those generators for which 
\begin{equation}\label{eq_EqualityProblem}
\mathcal{M}_{f, \lambda_1 , \dots , \lambda_m}^{[m]} 
\left( x_1 \,, \dots , x_m \right) 
= 
\mathcal{M}_{g, \lambda_1 , \dots , \lambda_m}^{[m]} 
\left( x_1 \,, \dots , x_m \right) 
\hspace{8mm} 
\left( x_1 \,, \dots , x_m \in C \right). 
\end{equation}
In the scalar valued setting when the generators 
are supposed the be continuous, 
then the solution is a well-known classical result: 
$ f = a \cdot g + b $ with some real numbers such that 
$ a \neq 0 $ (see \cite[Theorem 83 and Section 3.7]{HLP34} 
or \cite[Theorem A]{PP26}). 

In the case of non-continuous generators 
the situation is much more difficult. Even though it has been 
shown recently in \cite[Corollary 4.2]{PP26} 
that if the equality of 
two QAMs holds 
{\em for all} $ m \in \NN $, then the generators are 
affine transforms of each other, the same cannot 
be claimed in the case when we have 
{\em a fixed number of variables} $ m $. 
What is more, a fresh result of Kiss and Pasteczka 
gives examples for three-variable 
QAMs which are equal but their 
generators are not affine transforms of each other 
(see \cite[Example 1]{KP26}). 

The complete characterization of the solutions of the 
equality problem for fixed weights and 
a fixed number of variables is 
still open even in the real valued case. 
Therefore our aim here, instead of discussing 
the general equality problem, is to solve a particular 
case: for a fixed $ m \geq 2 $ when does the 
$ m $-variable vector-valued weighted quasi-arithmetic 
mean coincide with the weighted arithmetic mean? 
The $ m $-variable arithmetic mean (generated by 
the identity) with weights 
$ \lambda_1 \,, \dots , \lambda_m $ mean will be 
denoted by 
$ \mathcal{A}_{\lambda_1 , \dots , \lambda_m}^{[m]} $.  

\begin{prop}\label{prop-Application_invinj}
Let $ \emptyset \neq C \subseteq \Rn $ be a 
closed, convex set and let $ f : C \map \Rn $ 
be a strictly monotone mapping. Moreover, 
let $ m \geq 2 $ and $ \lambda_1 \,, \dots , \lambda_m > 0 $ 
such that 
$ \lambda_1 + \dots + \lambda_m = 1 $. 
Suppose that 
$ \mathcal{M}_{f, \lambda_1 , \dots , \lambda_m}^{[m]} 
= \mathcal{A}_{\lambda_1 , \dots , \lambda_m}^{[m]} $. 
Then, for all $ x \in C^{\circ} $, 
\[
y \in \conv (f(C)) \setminus \{ f(x) \} 
\ \mbox{ implies } \ 
f^{(-1)}(y) \neq x .
\]
%we have 
%$ \left( f^{(-1)} \right)^{-1} \left( f(x) \right) = 
%\lbrace f(x) \rbrace $. 
\end{prop}

\begin{proof} 
Firstly, $ \mathcal{M}_{f, \lambda_1 , \dots , \lambda_m}^{[m]} 
= \mathcal{A}_{\lambda_1 , \dots , \lambda_m}^{[m]} $ 
implies that, in particular, 
\begin{equation}\label{eq-prop_Appl_2mean}
\begin{aligned}
f^{(-1)} & \left( \lambda_1 f(x) + 
\left( 1-\lambda_1 \right) f(y) \right) = 
\mathcal{M}_{f, \lambda_1 , \dots , \lambda_m}^{[m]} 
\left( x,y, \dots , y \right) \\ 
& = \mathcal{A}_{\lambda_1 , \dots , \lambda_m}^{[m]} 
\left( x,y, \dots , y \right) = 
\lambda_1 x + \left( 1- \lambda_1 \right) y 
\hspace{15mm} 
\left( x,y \in C \right). 
\end{aligned} 
\end{equation}

Suppose that, contrary to the statement, 
there exists $ x \in C^{\circ} $ and 
$ y \in \conv(f(C)) $ such 
that $ f^{(-1)}(y) = x $ but $ y \neq f(x) $. 
Let us define the unit vector 
$ h := \frac{y-f(x)}{\nma y-f(x) \nma} $. 
Since $ x \in C^{\circ} $, there exists 
$ \alpha_0 > 0 $ such that 
$ x + \alpha h \in C $ for all 
$ \alpha \in (0,\alpha_0) $. Using the monotonicity 
of $ f^{(-1)} $ we obtain 
\begin{align*}
0 \leq  
\bsz	{f(x + \alpha h) - y}{ 
f^{(-1)} \left( f(x + \alpha h) \right) - f^{(-1)}(y) } 
= \alpha \cdot \bsz{f(x + \alpha h) - y}{h}, 
\end{align*}
so $ 0 \leq \bsz{f(x + \alpha h) - y}{h} $. 
%for every $ \alpha  \in (0,\tilde{\alpha}) $. 

%\begin{claim}
%According to the Gram-Schmidt orthogonalization, 
Now let us observe that, 
for all $ \alpha \in (0,\alpha_0) $, 
there uniquely exist 
$ \mu_{\alpha} > 0 $ and $ w_{\alpha} \in \Rn $ 
such that $ \mu_{\alpha} \geq \nma y - f(x) \nma $ 
and $ \bsz{w_{\alpha}}{h} = 0 $, moreover 
\begin{equation}\label{eq-prop_Appl-decomp}
f(x + \alpha h) = f(x) + \mu_{\alpha} h + w_{\alpha}. 
\end{equation}
%\end{claim}
%
%\begin{prclm}
Indeed, 
the decomposition \eqref{eq-prop_Appl-decomp} is 
due to the Gram--Schmidt algorithm with 
\[
\mu_{\alpha} = \bsz{f(x + \alpha h)-f(x)}{h} 
\hspace{3mm} \mbox{ and } \hspace{3mm} 
w_{\alpha} = f(x+ \alpha h) - f(x) - \mu_{\alpha} h, 
\]
%Using a previous observation we also have 
while the previous step of the proof yields 
\begin{align*}
0 & \leq \bsz{f(x + \alpha h) - y}{h} = 
\bsz{f(x + \alpha h) - f(x) + f(x) -y}{h} = 
\bsz{f(x+\alpha h)-f(x)}{h} \\ 
& + \bsz{f(x) - y}{h} 
= \mu_{\alpha} - \bsz{y-f(x)}{\frac{y-f(x)}{\nma y-f(x) \nma}} 
= \mu_{\alpha} - \nma y - f(x) \nma, 
\end{align*}
which is equivalent to 
$ \mu_{\alpha} \geq \nma y - f(x) \nma $. 
Therefore 
\begin{equation}\label{eq-prop_Appl_nu}
\nu := \inf 
\left \lbrace 
\bsz{f(x+\alpha h)-f(x)}{h} \ : \ 
\alpha \in \left( 0,\alpha_0 \right) 
\right \rbrace
\geq 
\nma y - f(x) \nma . 
\end{equation}
%\end{prclm}

In the next step let us choose 
%$ q \in \left( 1 , \frac{1}{1-\lambda_1} \right)$ and 
$ \beta \in \left( 0 , \alpha_0 \right) $ such that 
for $ z := x + \beta h \in C $ we have 
\begin{equation}\label{eq-prop_Appl-zchoice}
%z := x + \beta h \in C 
%\hspace{3mm} \mbox{ and } \hspace{3mm} 
\nu \leq \bsz{f(z)-f(x)}{h} < 
\frac{\nu}{\left( 1-\lambda_1 \right)}. 
\end{equation} 
Moreover, let us fix 
$ u := \frac{1+\lambda_1}{2} x + \frac{1-\lambda_1}{2} z 
= x + \frac{1-\lambda_1}{2} \beta h 
\in C $. Then 
\begin{align*}
0 & \leq 
\bsz{f^{(-1)} 
\bigl( \lambda_1 f(x) + (1-\lambda_1) f(z) \bigr) - 
f^{(-1)} \bigl( f(u) \bigr)}{ 
\bigl( \lambda_1 f(x) + (1-\lambda_1) f(z) \bigr) - f(u)} 
\\ 
& = 
\bsz{ \bigl( \lambda_1 x + (1-\lambda_1) z \bigr) - u}{ 
(1 - \lambda_1)f(z)-(1 - \lambda_1)f(x) + f(x) - f(u)} 
\\ 
& = 
\bsz{\frac{1-\lambda_1}{2} \beta \cdot h}{
\bigl( (1 - \lambda_1)(f(z)-f(x)) \bigr)
- \bigl( f(u) - f(x) \bigr)} 
\\ 
& = 
\frac{1-\lambda_1}{2} \beta \cdot 
\Bigl( (1-\lambda_1) 
\bsz{f(z)-f(x)}{h} - \bsz{f(u)-f(x)}{h} \Bigr) 
\\ 
& < 
\frac{1-\lambda_1}{2} \beta \cdot 
\Bigl( (1-\lambda_1) \frac{\nu}{1-\lambda_1} - \nu \Bigr) 
= 0, 
\end{align*}
where we have used \eqref{eq-prop_Appl_2mean}, 
\eqref{eq-prop_Appl_nu} and \eqref{eq-prop_Appl-zchoice} 
together with the monotonicity of $ f^{(-1)} $. 
However, the obtained contradiction means that 
$ f^{(-1)}(y) = x $ implies $ y = f(x) $, so the 
proof is complete by contraposition. 
\end{proof}


\begin{thm}\label{thm_Appl_vvQAMs}
Let $ m \geq 2 $, and let 
$ f : \Rn \map \Rn $ be a strictly increasing mapping. 
Moreover 
let $ m \geq 2 $ and 
$ \lambda_1 \,, \dots , \lambda_m > 0 $ 
such that 
$ \lambda_1 + \dots + \lambda_m = 1 $. 
Then the following two statements are equivalent: 
\begin{enumerate}
\item[\emph{(i)}] 
For all 
$ x_1 \,, \dots , x_m \in \Rn $ 
it holds that 
\[
\mathcal{M}_{f, \lambda_1 , \dots , \lambda_m}^{[m]} 
\left( x_1 \,, \dots , x_m \right)
= \mathcal{A}_{\lambda_1 , \dots , \lambda_m}^{[m]} 
\left( x_1 \,, \dots , x_m \right)
\] 
\item[\emph{(ii)}]
There exist a positive definite matrix 
$ A \in \mathcal{M}_{n \times n}(\RR) $ and a vector 
$ b \in \Rn $ such that 
\[
f(t) = At + b 
\hspace{8mm} 
\left( t \in \Rn \right). 
\]
\end{enumerate}
\end{thm}

\begin{proof} 
We begin with the proof of (i) $ \Longrightarrow $ (ii). 
From Proposition \ref{prop-Application_invinj} 
we have that $ f^{(-1)} $ is injective, so 
it coincides with the ordinary inverse. 
In particular, 
$ \conv \left( f \left( \Rn \right) \right) = 
f \left( \Rn \right) $, i.~e. $ f \left( \Rn \right) $ 
is convex. Suppose that $ t \in f \left( \Rn \right) $ 
is a boundary point of $ f \left( \Rn \right) $. 
Then there exist a supporting 
hyperplane for $ f \left( \Rn \right) $ at the 
point $ t $ (cf. \cite[Theorem 11.6]{Roc70}). 
This implies the existence of a vector 
$ 0 \neq h \in \Rn $ such that 
\[ 
\bsz{f(x)-t}{h} \geq 0 
\hspace{10mm} \left( x \in \Rn \right). 
\]
However, $ f $ is strictly increasing, so 
for $ x_0 := f^{-1}(t) $ we have 
\[ 
0< \bsz{f(x_0 - h)-f(x_0)}{
(x_0-h)-x_0} = 
\bsz{f(x_0 - h) - t}{-h}, 
\] 
which contradicts the choice of $ h $. Thus 
$ \partial \bigl( f \left( \Rn \right) \bigr) 
\cap f \left( \Rn \right) 
= \emptyset $, so $ f(\Rn) $ is open. 
Hence we may apply Brouwer's invariance of domain 
theorem for the injective, continuous 
(see Theorem \ref{thm-ExmoContinuous}) map 
$ f^{-1} : f \left( \Rn \right) \map \Rn $ and 
conclude that its inverse $ f : \Rn \to f(\Rn) $ 
is also continuous. 

Recalling 
\eqref{eq-prop_Appl_2mean} from the previous proof and 
applying $ f $ to both sides, we have 
\begin{equation}\label{eq-thm_Appl_lambda1}
\lambda_1 f(x) + (1-\lambda_1) f(y) 
= f \left( \lambda_1 x + (1-\lambda_1) y \right) 
\hspace{10mm} \left( x,y \in \Rn \right). 
\end{equation}
Consequently, every coordinate function 
$ f_i : \Rn \map \RR $ (for $ i = 1 , \dots , n $) of $ f $ 
%is $ \lambda_1 $-affine and continuous. It is 
%well-known that then each $ f_i $ 
is affine. That is, 
there exist $ a_i \in \Rn $ and $ b_i \in \RR $ 
such that 
\[
f_i(x) = \bsz{a_i}{x} + b_i 
\hspace{10mm} \left( i = 1, \dots , n \right)
\]
for all $ x \in \Rn $. The proof is similar to 
the solution process of Jensen's equation 
(which corresponds to \eqref{eq-thm_Appl_lambda1} with 
$ \lambda_1 = \frac{1}{2} $) discussed in 
\cite[Section 13.2]{Kuc09}. 
A slightly stronger statement can be found in 
\cite[Theorem 3.1]{BT26}. 
Summarizing the previous results, let 
$ A \in \mathcal{M}_{n \times n}(\RR) $ be the matrix 
which contains the vector $ a_i $ in its $ i $-th 
row and let the $ i $-th element of $ b \in \Rn $ 
be $ b_i $ (for $ i = 1 , \dots , n $). 
Then 
\[
f(x) = \left( f_1(x) , \dots , f_n(x) \right) = 
Ax + b 
\hspace{10mm} (x \in \Rn). 
\]
Furthermore, $ A $ has to be positive definite, because, 
for every $ 0 \neq x \in \Rn $, we have 
\[
\bsz{Ax}{x} = \bsz{(Ax+b)-(A \cdot 0+b)}{x-0} = 
\bsz{f(x)-f(0)}{x-0} > 0. 
\]

Finally, for the converse direction 
(ii) $ \Longrightarrow $ (i), 
let $ A $ be positive definite, $ b $ be arbitrary, 
and define $ f(t) = At+b $ for all $ t \in \Rn $. Then 
$ f^{-1}(y) = A^{-1}(y-b) $ for all $ y \in \Rn $, thus 
\begin{align*}
\mathcal{M}_{f, \lambda_1 , \dots , \lambda_m}^{[m]} 
& \left( x_1 \,, \dots , x_m \right)
= A^{-1} \bigl( \lambda_1 \left( A(x_1) + b \right) + 
\dots \lambda_m \left( A(x_m) + b \right) - b \bigr) \\  
& = 
A^{-1} \bigl( A \left( 
\lambda_1 x_1 + \dots + \lambda_m x_m \right) \bigr) = 
\mathcal{A}_{\lambda_1 , \dots , \lambda_m}^{[m]} 
\left( x_1 \,, \dots , x_m \right) 
\end{align*}
\end{proof}


\begin{thebibliography}{99}
\bibitem[BT26]{BT26} Z.~Boros, P.~Tóth, {\em Weakly affine functions}, Publ. Math. Debr., {\bf 108}/3-4 (2026), 295--326.
\bibitem[HLP34]{HLP34} G.~H.~Hardy, J.~E.~Littlewood, G.~P\'olya, {\em Inequalities}. University Press, Cambridge, 1934.
\bibitem[KP26]{KP26} T.~Kiss, P.~Pasteczka, {\em Equality problem for generalized quasiarithmetic means generated by discontinuous strictly monotonic functions}, (2026), https://doi.org/10.48550/arXiv.2607.22266
\bibitem[Kuc09]{Kuc09} M.~Kuczma, {\em An Introduction to the Theory of Functional Equations and Inequalities}. Birkhäuser Basel, 2009.
\bibitem[Leo23]{Leo23} P.~Leonetti, {\em Quasi-Arithmetic Means Ad Libitum}, Results Math. {\bf 78}:238 (2023)
\bibitem[PP26]{PP26} Zs.~Páles, P.~Pasteczka, {\em Equality and comparison of generalized quasi-arithmetic means}, Rev. Real Acad. Cienc. Exactas Fis. Nat. Ser. A-Mat., {\bf 120}:19 (2026)
\bibitem[Roc70]{Roc70} R.~T.~Rockafellar, {\em Convex Analysis}, Princeton University Press, 1970.
\end{thebibliography}
\end{document}
